\documentclass[11pt]{article}

\usepackage[utf8]{inputenc}
\usepackage[T1]{fontenc}
\usepackage[margin=1in]{geometry}
\usepackage{graphicx}
\usepackage{booktabs,longtable,tabularx,array,multirow}
\usepackage{ragged2e}
\usepackage{pdflscape}
\usepackage{float}
\usepackage{placeins}
\usepackage{caption}
\usepackage{rotating}
\usepackage{makecell}
\usepackage[table]{xcolor}
\usepackage{xcolor}
\usepackage{amssymb}
\usepackage{microtype}
\usepackage{csquotes}
\usepackage{xurl}
\usepackage[hidelinks]{hyperref}


\usepackage{fontspec}
\usepackage{xeCJK}

\setmainfont{TeX Gyre Termes}
\setCJKmainfont{Noto Serif CJK SC}

% 日文单独指定字体
\newCJKfontfamily\japanesefont{Noto Serif CJK JP}


\usepackage[
    backend=biber,
    style=numeric-comp,
    sorting=none,
    sortcites=true,
    natbib=true
]{biblatex}

\addbibresource{references.bib}
\title{Toward Objective Stroke Assessment: An Item-Level Analysis of Clinical Rating Scales for AI-Based Motion Analysis}

\author{
    Siqi Yang$^{1}$ \and
    Wenbo Qi$^{1}$ \and
    xxx$^{3}$ \and
    xxx$^{3}$ \and
    Baosheng Yu$^{1,2,3}$
}
\date{
    $^{1}$Lee Kong Chian School of Medicine, Nanyang Technological University, Singapore \\
    $^{2}$Centre of AI in Medicine, Nanyang Technological University, Singapore\\
    $^{3}$Rehabilitation Research Institute of Singapore, Nanyang Technological University, Singapore \\[1ex]
    \texttt{\{xxx,xxx\}@ntu.edu.sg, xxx@xxx.org}
}

\begin{document}
\maketitle

\begin{abstract}
Stroke is a leading cause of long-term disability worldwide, making objective assessment essential for consistent evaluation of motor function, personalized rehabilitation, and longitudinal recovery monitoring. Conventional clinical rating scales remain subjective and rely heavily on expert clinicians, limiting their consistency, scalability, and accessibility, particularly for frequent and remote assessments. 
Recent advances in computer vision, wearable sensing, and machine learning have enabled AI-based motion analysis for automated stroke assessment. However, most existing approaches formulate stroke assessment as end-to-end prediction of clinical scores, thereby inheriting the subjectivity of the underlying rating scales while providing limited insight into the quantitative movement features that underpin clinical assessment.
This survey presents a component-level framework that links clinical assessment with AI-based motion analysis. Rather than treating clinical scores as prediction targets, it examines how established rating scales can be decomposed into quantitative movement features that AI can objectively quantify. The survey reviews clinical rating scales, maps their assessment criteria to measurable movement features, and discusses sensing modalities and AI techniques for objective movement analysis. By explicitly linking clinical assessment with quantitative movement analysis, it provides a more interpretable and clinically grounded framework for AI-based stroke assessment and a foundation for developing next-generation objective clinical rating scales.
\end{abstract}


\section{Introduction}

Stroke is a leading cause of long-term disability worldwide, making accurate and objective assessment essential for personalized rehabilitation and longitudinal recovery monitoring. Across the stroke continuum, human movement serves as a key functional biomarker, reflecting both pre-stroke risk and post-stroke impairment and recovery. Consequently, objective quantification of human movement provides a natural foundation for stroke assessment. However, current clinical practice relies predominantly on standardized rating scales based on clinician observation and qualitative judgment rather than quantitative analysis of movement. Widely used scales, such as the Fugl-Meyer Assessment (FMA), summarize motor performance through coarse ordinal ratings that are inherently subjective, susceptible to inter-rater variability, and dependent on trained clinicians, limiting their scalability and suitability for frequent or continuous assessment.

Recent advances in computer vision, wearable sensing, and machine learning have enabled AI-based motion analysis to quantify human movement from video and wearable sensor data. Most existing approaches formulate stroke assessment as end-to-end prediction of clinical scores, learning direct mappings from raw motion signals to clinician-assigned ratings without explicitly modeling the quantitative movement features that underpin clinical assessment. Consequently, these models inherit the subjectivity and limited granularity of the target rating scales while providing little insight into how specific movement abnormalities contribute to clinical outcomes. More fundamentally, treating clinical scores as the primary prediction target overlooks the fact that these scores are higher-level clinical interpretations of underlying movement characteristics. As illustrated in Figure~\ref{fig:figure1}, effective AI-based stroke assessment should therefore explicitly quantify clinically meaningful movement features and relate them to clinical evaluation, rather than directly predicting clinical scores alone. This perspective enables more interpretable representations of motor dysfunction and better supports rehabilitation planning, longitudinal monitoring, and clinical decision-making.

\begin{figure}[!t]
    \centering
    \includegraphics[width=0.9\linewidth,height=0.55\linewidth]{figure1.pdf}
    \caption{Overview of the component-level framework linking clinical rating scales with AI-based motion analysis. Instead of directly predicting clinical scores, the framework decomposes clinical assessment into quantitative movement features that can be objectively measured by AI, enabling more interpretable and clinically grounded stroke assessment.}
    \label{fig:figure1}
\end{figure}


Motivated by this perspective, this survey introduces a component-level framework that bridges established clinical assessment with AI-based motion analysis. Rather than treating clinical scores as end-to-end prediction targets, it investigates how widely used stroke assessment scales can be decomposed into quantitative movement features that AI can objectively quantify. Specifically, the survey reviews established clinical rating scales, maps their assessment criteria to quantitative movement features, and systematically examines AI-based motion analysis methods through this lens. It further evaluates the extent to which current approaches capture clinically meaningful aspects of motor function and identifies important gaps between clinical assessment requirements and existing computational capabilities. By explicitly connecting clinical evaluation to quantitative movement analysis, this survey provides a roadmap for developing more interpretable, clinically grounded, and objective AI systems for stroke assessment.

The remainder of this paper is organized as follows. Section~\ref{sec:background} provides the background on stroke assessment and AI-based motion analysis. Section~\ref{sec:scales} reviews commonly used clinical rating scales and analyzes their constituent assessment components from a motion analysis perspective. Section~\ref{sec:framework} presents the proposed component-level framework for objective AI-driven stroke assessment. Finally, Section~\ref{sec:challenge} discusses the key challenges, emerging opportunities, and future research directions toward objective and interpretable AI-driven stroke assessment.


\section{Background and Preliminaries}
\label{sec:background}

Stroke results from cerebral ischaemia or haemorrhage; motor impairment disrupts voluntary control, coordination, balance, gait and upper-limb function, restricting mobility and daily activities \cite{WHO2025StrokeFactSheet,S1VADoD2024StrokeRehab}. Assessment recurs along the care pathway with changing purpose: in acute care, the National Institutes of Health Stroke Scale (NIHSS) grades neurological severity and early change \cite{S1Prabhakaran2026AHAASAAIS,ScaleBrott1989NIHSS}; in rehabilitation and after discharge, assessment covers motor impairment, balance, gait, upper-limb activity and independence to support goal setting and follow-up \cite{S1Winstein2016AHAASARehab,S1VADoD2024StrokeRehab,S1NICE2023NG236}. The resulting ordinal ratings, times, speeds, distances and counts describe complementary aspects of recovery and are not interchangeable \cite{ScaleHarrison2013Clinimetric,CoreKwakkel2017SRRRSensorimotor,CoreVanCriekinge2024SRRRBalanceMobility}.

We distinguish four levels of description: the \emph{clinical construct} a scale measures (for example, standing balance); the \emph{score-defining criterion} that separates adjacent grades (standing up with or without the hands); the \emph{observable movement information} in that criterion (hand use, duration, loss of balance); and \emph{computational features} derived from motion signals (gait speed, joint angle, trunk displacement). A criterion may combine observable movement with information that cannot be observed, and a feature may quantify a criterion, indicate it indirectly or have no counterpart in any scale.

Pose estimation recovers body landmarks and trajectories from RGB or depth video, giving the timing, displacement and joint angles of visible segments with an accuracy that depends on viewpoint, calibration, occlusion and tracking \cite{S3Luvizutto202234115576,S3Lonini202235224426,S3John202437682075}. Inertial measurement units (IMUs) record segment acceleration and angular velocity, also outside a camera's view \cite{S3OBrien202236304845,Yang2013SpatiotemporalGait}. Neither records force, muscle tone or what the patient perceives. These signals yield five feature families that describe how a movement is executed (Table~\ref{tab:computational_feature_families}).

\section{Clinical Rating Scales: A Motion Analysis Perspective}
\label{sec:scales}

\subsection{Review Methodology}

\paragraph{Tool library.} We assembled a purposive library of standardized motor and mobility assessments for criterion-level mapping, not an exhaustive inventory (Fig.~\ref{fig:figure2}). A guideline stream drew on adult-stroke guidelines, standards and consensus documents that make assessment recommendations and have a complete English full text, identified in PubMed, Embase, the Guidelines International Network, TRIP, World Health Organization resources, issuing-body websites and reference lists (searches completed May 2026). Two reviewers screened records independently and resolved disagreements by consensus. Of 41 full-text documents, 13 were excluded, seven for lacking a complete English text; the 28 included guidelines from nine regions named 37 tools (Supplementary Table~\ref{tab:s1_guidelines_compact}). A supplementary stream added standardized tools used in stroke research but absent from these guidelines: of 240 candidates from PubMed-indexed stroke studies published 2022--2026 and from citation tracking, 184 were supported by at least one meta-analysis, systematic review or multicentre study, and 22 were retained. The Modified Tardieu Scale, recommended only in an excluded non-English guideline, was added from its English specification \cite{S2RehabMeasuresModifiedTardieu}. Merging walking speed into the 10-Metre Walk Test (10MWT) and removing two records without standardized motor scoring (observational gait analysis and the Frenchay Activities Index) left 57 tools. Search strings and screening logs are in Supplementary Data~1.

\paragraph{COA classification.} Tools were classified by how the score is generated, following clinical outcome assessment (COA) definitions \cite{FDANIH2016BESTCOA,Walton2015COAConceptualFoundation}: performance outcomes (PerfO) when grades are measured quantities or observable task outcomes, and clinician-reported outcomes (ClinRO) when grades rest on clinician judgement of assistance, supervision, safety, movement normality, examination signs, recovery stage or global status, even for standardized tasks. The library comprised 16 PerfO, 33 ClinRO and six patient-reported outcome (PRO) tools and two tools with mixed outputs: the Wolf Motor Function Test (WMFT; timed and strength outputs PerfO, Functional Ability Scale ClinRO) and the Rivermead Mobility Index (RMI; 14 self-reported items and one observed item) (Supplementary Table~\ref{tab:supplementary_tool_library_clinical}). The classification is operational, not regulatory.

\paragraph{Mapping entries.} Each tool was mapped at the level at which its scoring criteria form distinct rules, giving 196 entries. Thirteen tools with item-specific ordinal criteria, such as the Berg Balance Scale (BBS), Trunk Impairment Scale (TIS), Mini-Balance Evaluation Systems Test (Mini-BESTest) and Tinetti Performance-Oriented Mobility Assessment (POMA), were mapped item by item (129 entries); ten tools whose items share criteria within subtests, such as the Action Research Arm Test (ARAT), Fugl-Meyer Assessment (FMA) and WMFT, by subtest (33 entries); and 34 tools as a whole (34 entries), mostly tests with a single time, distance, count or category output, such as the 10MWT, Timed Up and Go (TUG) and Functional Ambulation Category (FAC), global ratings and questionnaires. Some item-structured tools (Function in Sitting Test, Community Balance and Mobility Scale, Chedoke Arm and Hand Activity Inventory) were mapped as whole tools, so tool-weighted proportions are reported alongside entry counts. For multi-domain tools only motor and mobility content was mapped, following the versions in Supplementary Table~\ref{tab:supplementary_tool_library_clinical}.

\paragraph{Coded attributes.} Four attributes were coded for each entry from its official scoring criteria, including general rules of the tool that change an item's score; rules and worked examples are in Supplementary Data~1. (i) \emph{Movement components} were marked only when the criteria grade performance in that component, not when the task merely requires it. (ii) The \emph{information type} distinguishing grades was recorded (kinematics or movement quality; time, distance or count; assistance or supervision; safety judgement; verbal cue or response; applied force or touch; tone or reflex; self-report; more than one allowed). (iii) \emph{Availability from ordinary video} (one RGB camera framing the task as administered) was Full when all grade-defining information is visible under adequate recording, including assistance or support that occurs; the presence of an examiner, guarding, standard instructions or visible physical help did not by itself lower the code. It was Partial when some of that information is not visible: the amount of assistance rather than its occurrence, examiner-applied force or resistance, events confirmed by touch, verbal cues or responses, judgements without a visible trace, or usual performance outside the recorded occasion. It was None when the score does not come from visible performance (patient report, reflexes, or tone and force read by touch or instrument). For example, BBS standing items graded partly by supervision were Full because the instability that calls for supervision is visible, whereas BBS sit-to-stand was Partial because its lower grades distinguish minimal from moderate or maximal assistance; ratings of observed or usual performance by a clinician, as in the Barthel Index (BI), Functional Independence Measure (FIM) and modified Rankin Scale, were Partial, whereas RMI questionnaire items, answered by the patient, were None. None refers to video alone: such content can be obtained by interview and combined with video-derived measures. (iv) Each entry was linked to the \emph{computational feature families} that can quantify the observable part of its criteria (Table~\ref{tab:computational_feature_families}). Availability denotes criterion-level observability, not measurement accuracy or automated scoring, and family links denote classes of measures, not validated features. Supplementary Data~1 gives, for every entry, the task, the full scoring criteria, the rationale for its code, the information video cannot provide with an alternative source, and the recording requirements.

\paragraph{Verification.} A second reviewer checked the availability codes of the initial coding (193 entries) with the first reviewer's codes visible. Agreement was 77.2\% (Cohen's $\kappa$ = 0.63; linear-weighted $\kappa$ = 0.68 with Full, Partial and None ordered); 40 of the 44 disagreements were Partial versus Full, and 34 entries were recoded after checking the scoring manuals. The final 196 entries comprise 178 checked entries that were unchanged, three that were recoded later and 15 that were added or redefined after the check. The availability rule was then revised so that visible assistance no longer lowered the code by itself, and one coder re-examined all 196 entries against the scoring manuals, changing 36 codes and 19 information types and flagging 22 boundary judgements. Movement components, information types and feature-family links were assigned by the same coder. The agreement statistics therefore apply only to the initial coding under the earlier rule (Section~\ref{sec:limitations}); every change and its source is listed in Supplementary Data~1.

\begin{figure}[!htbp]
    \centering
    \includegraphics[
        height=0.58\textheight,
        width=\linewidth,
        keepaspectratio
    ]{figure2.pdf}
    \captionsetup{justification=raggedright,singlelinecheck=false}
    \caption{\textbf{Construction of the assessment-tool library.} The diagram summarizes library construction, not a record-level PRISMA flow. \textbf{a,} Guideline stream; record counts are given for the PubMed cross-check only, and full-text assessment pooled all sources. \textbf{b,} Supplementary stream; per-reason counts for leads not retained were not tabulated. Harmonization merged walking speed into the 10MWT and removed two records without standardized motor scoring. PerfO, ClinRO and PRO, performance, clinician-reported and patient-reported outcomes (Section~\ref{sec:scales}); GIN, Guidelines International Network; TRIP, TRIP Database; WHO, World Health Organization.}
    \label{fig:figure2}
\end{figure}


\subsection{Clinical Rating Scales}
\label{sec:clinical_rating_scales}

The main text synthesises 13 representative tools (Table~\ref{tab:representative_scales_compact}): 11 reflect priorities of the three Stroke Recovery and Rehabilitation Roundtable (SRRR) consensus rounds \cite{CoreKwakkel2017SRRRSensorimotor,CoreKwakkel2019SRRRUpperLimb,CoreVanCriekinge2024SRRRBalanceMobility}, and the TUG and BI were added from guidelines for transitional mobility and daily activities \cite{S1Winstein2016AHAASARehab,S1VADoD2024StrokeRehab}. Together they cover the three kinds of score output in the library: ordinal item criteria; time, speed or distance; and categorical independence ratings.

Recurrent item content was grouped into six movement components: joint mobility and range of motion; muscle strength and force production; coordination and selective motor control; balance and postural control; manual dexterity and upper-limb task execution; and gait and functional mobility. Movement quality and compensation are treated as modifiers across components, because the same task outcome can be reached through restitution or compensation \cite{Levin2009RecoveryCompensation,CoreKwakkel2019SRRRUpperLimb}. The grouping follows the International Classification of Functioning, Disability and Health (ICF) and SRRR domains \cite{WHO2001ICF,Cieza2019ICFLinking,CoreKwakkel2017SRRRSensorimotor} and supports cross-scale comparison; it is not a validated taxonomy. A component is marked in Table~\ref{tab:representative_scales_compact} only when the criteria grade it: walking 10~m requires balance, but the 10MWT records speed, so only gait is score-defining; postural control determines most BBS and Mini-BESTest grades but is not rated in the ARAT.

Shared components do not make scores interchangeable. The FMA upper-extremity motor subscale (FMA-UE) grades synergy-dependent and selective movement \cite{ScaleFuglMeyer1975PhysicalPerformance,ScaleSee2013StandardizedFMA}, whereas the ARAT grades completion of 19 tasks \cite{ScaleYozbatiran2008StandardizedARAT}; a task completed through trunk or shoulder compensation can score on an activity measure without restored motor control \cite{Levin2009RecoveryCompensation,Santisteban2016UpperLimbOutcome}. The TIS, BBS and Mini-BESTest all grade postural control but differ in body segment, task context and scoring information \cite{ScaleVerheyden2004TIS,ScaleBerg1992BalanceValidation,ScaleFranchignoni2010MiniBESTest}. Lower-limb tools range from impairment (FMA lower-extremity subscale, FMA-LE) through walking independence (FAC) and speed and endurance (10MWT; 6-Minute Walk Test, 6MWT) to adaptive and transitional mobility (Dynamic Gait Index, DGI; TUG) \cite{ScaleHolden1984FAC,ScaleCheng2020StrokeWalkTests,ScaleAn2017DGIStroke,ScaleNg2005TUGStroke}; because SRRR considers the 10MWT, 6MWT and DGI untestable when the FAC is below 3 \cite{CoreVanCriekinge2024SRRRBalanceMobility}, the observable content depends on walking independence. Only the motor and mobility items of the NIHSS and BI relate to movement \cite{ScaleBrott1989NIHSS,ScaleMahoney1965Barthel}.

\input{Table1}
\clearpage

\section{Toward Objective AI-Driven Stroke Assessment}
\label{sec:framework}
\subsection{From Clinical Scores to Computational Feature Families}

AI models can estimate selected scores or categories from motion data \cite{Lee2021SmartphoneVideoAmbulation,Formstone2021IMUMMG,Zamin2023BIONICS,Li2017IMUEMGMotorEvaluation}, in some studies item by item; for example, 13 FMA-UE items were scored from depth-camera recordings with per-item accuracies of 65--87\% \cite{Kim2016KinectFMA}. Yet the same score can arise from restricted excursion, slow execution, asymmetry, instability or compensation, which imply different rehabilitation priorities, and SRRR recommends kinematic and kinetic measures alongside clinical scales \cite{CoreKwakkel2019SRRRUpperLimb,CoreVanCriekinge2024SRRRBalanceMobility}. Our mapping asks which grades are defined by information that a recording can supply, following the chain clinical item $\rightarrow$ scoring criterion $\rightarrow$ observable information $\rightarrow$ computational feature $\rightarrow$ clinical interpretation. In the BBS sit-to-stand item, hand push-off, time to stand and loss of balance are observable as wrist--armrest contact, transition duration and trunk sway, but the lower grades are defined by how much assistance the examiner gave, which no motion feature encodes. In the TUG, detecting onset and the return to sitting measures the score itself, and phase durations and turning steps add information the score lacks.

Components and feature families relate many-to-many: a gait item yields temporal, spatial and relational measures, and one joint-angle computation can quantify FMA range of motion or describe trunk inclination during TIS items. Features also differ in how closely they correspond to a criterion (Table~\ref{tab:computational_feature_families}). Some quantify the criterion itself, such as test time, walking speed, reach distance or active joint excursion where range is graded. Others only indicate it: joint coupling or trunk displacement may reflect abnormal synergy or compensation, but the criterion is a clinical judgement of the pattern, so an accurately computed feature may still fail to represent the construct. Measurement adds further limits: distance and joint angles need calibration, an aligned view and reliable tracking, and force, passive tone, reflexes, sensation, pain and self-report cannot be read from video, although parts of an examination, such as the angle of a catch during passive stretch, may be visible.

\input{Table2}

Across the 196 entries, score-defining information was fully available from video for 134 (68\%), partially for 36 (18\%) and not at all for 26 (13\%) (Fig.~\ref{fig:video_availability}a). These proportions describe the information content of scoring criteria, not how much scoring can be automated. The 104 entries graded only by movement, time, distance or count were Full by construction; the informative result concerns the others. Of 55 entries in 19 tools that combined movement with assistance, supervision, safety or verbal information, 30 were Full because the defining assistance, support or instability is a visible event, and 25 were Partial because grades depend on the amount of assistance (nine), verbal cues or responses (eight), usual performance outside the recording (six) or a supervision judgement without a visible trace (two). Entries graded partly by examiner-applied force, tone or reflexes (17; 11 Partial, six None) or by patient report (20) were never fully available.

The distribution depends on the unit of analysis, because item-level balance scales such as the TIS, POMA, Mini-BESTest and BBS contribute many entries (Fig.~\ref{fig:video_availability}d). Weighting tools equally lowered the fully available share to 53\% and raised the unavailable share to 18\% (Fig.~\ref{fig:video_availability}b). The Rivermead Mobility Index provided 15 entries (8\%) but 14 of the 26 unavailable ones; without it, 133 of 181 entries (73\%) were fully, 36 (20\%) partially and 12 (7\%) not available. Among the 176 entries based on observed performance or examination, 134 (76\%) were fully, 36 (20\%) partially and six (3\%) not available, all six graded by force, tone or reflexes. Under every weighting, the limits of video lie in examination, self-report, verbal and usual-performance content rather than in the visibility of movement or assistance.

Feature-family anchoring was similarly uneven: temporal and event measures were anchored in 153 entries from 42 tools, spatial and trajectory measures in 81 from 24, joint and segment kinematics in 52 from 21, inter-segment and relational measures in 46 from 19, and trajectory-derived quality in six from six (Fig.~\ref{fig:video_availability}c). Scales designed for visual rating thus provide many anchors for timing and kinematic features but few for the smoothness and jerk measures prominent in AI studies, which does not make movement quality clinically unimportant (Section~\ref{sec:constructs}).

\begin{figure}[!htbp]
    \centering
    \includegraphics[width=\linewidth]{figure3.pdf}
    \caption{\textbf{What ordinary video can provide for scale scoring.} \textbf{a,} Information profile of each of the 196 mapping entries, by what its scoring criteria use (left), and its availability from ordinary video (right); band width is proportional to entry count. Entries graded only by movement, time, distance or count are Full by definition. \textbf{b,} Availability weighted by entry or by tool. \textbf{c,} Entries and tools whose criteria anchor each feature family; an entry can anchor several families. \textbf{d,} One square per entry for the five tools with the most entries in each domain (all tools in Supplementary Data~1). Full, all score-defining information visible under an adequate recording, including assistance that occurs; Partial, some grade-defining information not visible (amount of assistance, verbal cues, force, palpation, judgements without a visible trace or usual performance); None, score not derived from visible performance. Availability refers to whether a rater could assign the grade from an adequate recording, not to automated accuracy. Percentages may not sum to 100 because of rounding. Abbreviations follow Supplementary Table~\ref{tab:supplementary_tool_library_motion}.}
    \label{fig:video_availability}
\end{figure}
\clearpage

\subsection{Illustrative Examples Across Feature Families}
The following studies, identified during mapping rather than by systematic search (see dedicated reviews \cite{S3Luvizutto202234115576,Senadheera2024AIStrokeScoping,Nicora2025AIRoboticRehab}), illustrate what each feature family has measured and how it was validated.

\emph{Temporal and event measures} map most closely onto scoring criteria. IMUs and single-camera pose estimation have measured walking speed, cadence and gait-phase timing against instrumented references, and joint and object tracking has supported upper-limb action segmentation and task-completion detection \cite{Yang2013SpatiotemporalGait,S3Lonini202235224426,S3Stenum202438530801,Ahmed2021AutomatedMovement}; their validity rests on reproducible definitions of onset, termination and phase transitions. \emph{Spatial and trajectory measures} such as reaching paths, endpoint behaviour, sway and gait paths have been obtained with depth and RGB systems \cite{Scano2018KinectReaching,Proffitt2023DepthMovement,Zamin2023BIONICS,S3Stenum202438530801}; metric distances require calibration, and viewpoint, occlusion and keypoint failure distort paths. \emph{Joint and segment kinematics} have been measured with IMUs, exoskeleton sensors and video pose estimation against motion capture \cite{Held2018UpperLimbIMU,Grimm2021KinematicValidation,Formstone2021IMUMMG,S3John202437682075}; such agreement shows that an angle is measured correctly, not that it captures the synergy or selective control the FMA grades.

\emph{Inter-segment and relational measures} (symmetry, phase, joint coupling, trunk--limb relations) from smartphone video, Kinect and multi-IMU systems have been used to classify ambulation dependence, compensation, fall risk or later walking function \cite{Lee2021SmartphoneVideoAmbulation,Zhi2018CompensationDetection,S3Abdollahi202438339529,S3OBrien202236304845}; high accuracy shows that a labelled pattern is detectable, not that the features measure coordination, compensation or risk. \emph{Trajectory-derived quality} measures (smoothness, jerk, movement units, path efficiency, variability) have been used to classify therapist-rated quality, monitor rehabilitation or separate correct from compensatory movements \cite{Jung2017MovementQuality,Lu2021MotionFeaturesProgress,Proffitt2023DepthMovement,Lin2019ManualDexterity,Seo2024WearableQuality}, in one case with a wrist-IMU model trained on simulated movement in neurotypical adults \cite{Li2023DeepLearningTaskPractice}; because few scoring criteria grade these qualities, such studies rely on therapist labels or study-specific references.

\subsection{Integrating Feature-Level Measures for Clinical Interpretation}
A total score does not show how it was achieved. Two patients with the same TUG time may lose time in rising or in turning, and two may reach a target through elbow extension or through trunk flexion; these differences point to different training targets. Yet a single feature rarely determines a clinical state: trunk displacement during reaching may reflect compensation, task strategy or seating, and its meaning depends on joint excursion, endpoint accuracy and task completion. Feature-level measures should therefore be reported as a task-specific profile rather than recombined into another unvalidated total. Each measure should state whether it quantifies a criterion or only indicates a construct (Table~\ref{tab:computational_feature_families}) and report the protocol, assistance context, uncertainty and a normative or within-person reference; non-motor information should be recorded separately, and established scales remain the reference.

Such profiles require staged evidence. Adapting the V3 framework for digital measures, which separates verification, analytical validation and clinical validation \cite{Goldsack2020V3}, and adding clinical utility, we distinguish \emph{analytical validity} (a feature is measured accurately and reliably against a reference), \emph{task-level validity} (repeatability and responsiveness under a defined protocol), \emph{clinical-construct validity} (the measure represents the intended impairment or activity) and \emph{clinical utility} (its use improves assessment or decisions). Pose error, test--retest reliability, classification accuracy and correlation with a scale address different levels and are not interchangeable endpoints (Section~\ref{sec:validation}); measurement properties follow the COSMIN taxonomy \cite{Mokkink2010COSMINTaxonomy}.

\section{Challenges and Future Directions}
\label{sec:challenge}

\subsection{Standardisation of Quantitative Movement Features}

Even temporal and event measures, the features most often anchored in scoring criteria, depend on how onset, termination and phase transitions are defined. For relational and quality features, one name can hide different equations, event rules, reference frames and normalisations, and values change with sampling rate, filtering and differentiation \cite{Jung2017MovementQuality,Proffitt2023DepthMovement,Yang2013SpatiotemporalGait,S3Stenum202438530801}; a systematic review of upper-limb kinematics after stroke found 151 metrics in 225 studies, of which 62 had any evaluation of measurement properties \cite{Schwarz2019UpperLimbKinematicsSR}. Feature definitions should be reported as part of the protocol (task and instructions, coordinate system and units, event rules, sampling rate, filtering, missing-data handling, formula and trial aggregation), and benchmarks pairing raw signals with annotated events on clinically defined tasks would separate analytical choices from patient performance. Standardisation improves comparability but does not establish clinical meaning.

\subsection{From Computational Features to Clinical Constructs}
\label{sec:constructs}

Asymmetry, joint coupling, trunk displacement and smoothness can be measured, but each may reflect impairment, adaptive strategy, task constraints, assistance or measurement error; classifying therapist labels or correlating features with scale scores gives concurrent evidence, not construct equivalence \cite{Zhi2018CompensationDetection,Lee2021SmartphoneVideoAmbulation,Li2017IMUEMGMotorEvaluation,Grimm2021KinematicValidation}. The sparse anchoring of movement quality (six entries) has two readings: features without a criterion-level anchor may lack established clinical meaning, or visually rated scales may omit dimensions that instrumented measurement can add, as SRRR recognises \cite{CoreKwakkel2019SRRRUpperLimb}. The two cases need different validation. A feature meant to reproduce or explain a grade should be validated against that item; a feature meant to capture what the scale misses, such as residual deficits near the ceiling or restitution versus compensation, cannot be validated by agreement with the scale and needs known-group, longitudinal and independent-measure evidence. Either way, studies should pre-specify the component each feature represents and whether the link is direct or proxy, validate at task and item level before estimating totals, test across severity, assistance, devices and settings, and use multi-rater labels with reported reliability, because rater grades are themselves imperfect references.

\subsection{Score-Defining Content Beyond Movement Kinematics}

More than a quarter of entries (55 of 196) were graded partly by assistance, supervision, safety or verbal information. In 30 of them the grade is defined by visible events, such as therapist contact, use of an aid or support surface, loss of balance or corrective steps, which can be annotated and detected alongside kinematics; a system that ignores these events cannot reproduce how the items are scored. The other 25 depend on information the image does not carry: the share of effort a helper provided (Ability for Basic Movement Scale, Function in Sitting Test), verbal cues or responses (BBS transfer and turning items, the Mini-BESTest dual-task item, the supervision category of the FAC) and supervision decisions without a visible event. These should be recorded by the clinician or captured with audio rather than inferred from motion. Daily-activity tools such as the BI and FIM also rate usual performance, whereas a recording shows one occasion. Datasets should therefore annotate assistance events, the amount of help, cues and context separately, so that models are neither credited with nor penalised for information the recording lacks.

\subsection{Multimodal Assessment Beyond Video}

Video cannot supply examiner-applied force and resistance, passive tone and reflexes (17 entries), sensation, pain and patient-reported difficulty (20 entries), verbal content or movement outside the camera's view \cite{S3John202437682075,S3Stenum202438530801}. IMUs add segment orientation and acceleration outside a fixed view \cite{Yang2013SpatiotemporalGait,Held2018UpperLimbIMU}; electromyography and mechanomyography provide muscle-related signals; and pressure, force and robotic sensors quantify contact loads and interaction dynamics \cite{Li2017IMUEMGMotorEvaluation,Formstone2021IMUMMG,Grimm2021KinematicValidation}. These signals are not the clinical quantities themselves: muscle activity is not strength, and device-measured resistance is not a Modified Ashworth grade without validation. Sensors should be added when a predefined question needs information that video lacks, such as pressure sensing for weight transfer in sit-to-stand or force measurement where a criterion grades strength, with synchronised clocks and validated inputs. Because each sensor adds setup time, cost and burden, a tiered design, with video for repeated measurement and instrumented sensing only where force, muscle activity or out-of-view motion is essential, may balance completeness against feasibility, acceptability and privacy.

\subsection{Validation and Clinical Utility}
\label{sec:validation}

Different evidence supports different conclusions. Analytical validity shows that a quantity is measured correctly; reliability, that repeated measurements agree; construct validity, from convergent, discriminant and known-group comparisons, that a measure reflects the intended impairment or activity, a correlation with a scale being only part of this evidence; responsiveness, that meaningful within-person change is detected; external generalisability, that performance holds across patients, sites, devices and recording conditions; and clinical utility, that use of the measure improves assessment, decisions or outcomes. The examples cited here are mostly proof-of-concept cohorts, internally cross-validated models or correlations with clinical scores \cite{S3Lonini202235224426,Lee2021SmartphoneVideoAmbulation,Proffitt2023DepthMovement,Grimm2021KinematicValidation}, and a recent scoping review of AI for spasticity assessment, mostly in post-stroke upper limbs, found no externally validated study \cite{Cao2026AISpasticity}; such designs establish feasibility or association, not generalisability or usefulness. Future evaluations should pre-specify the intended use, population, task and comparator, test prospectively across severity, assistance levels, devices and sites with subgroup performance and failure rates, and compare utility with standard assessment, including setup time, workflow integration, acceptability and cost.

\subsection{Limitations}
\label{sec:limitations}

Five limitations apply. First, the library was built purposively from English-language guidelines and supplementary literature, not by exhaustive search; tools used mainly where guidelines were excluded, such as Japan, or lacking an English specification may be missing, and per-reason exclusions in the supplementary stream were not tabulated. The proportions describe widely recommended instruments, not stroke assessment in general.

Second, the mapping unit was not set by a single a priori rule, and entry counts favour item-level balance and mobility scales. The unavailable share ranged from 13\% with entry weighting to 18\% with tool weighting and fell to 7\% without the Rivermead Mobility Index; the pattern was consistent, but percentages should not be compared with analyses that use another unit.

Third, coding was only partly verified. The availability check was not blinded, and the moderate agreement ($\kappa$ = 0.63; weighted $\kappa$ = 0.68) shows that the Full--Partial boundary requires judgement. The re-examination under the revised rule, which changed 36 codes, was not double coded; 22 entries rest on boundary judgements, such as whether visible instability shows the need for supervision or whether an optional verbal cue changes a grade; and movement components, information types, feature-family links and 15 new or redefined entries were coded by one person. Results resting on them, particularly the feature-family distribution and the sparse anchoring of movement quality, are provisional until independent double coding is available.

Fourth, availability was judged from written criteria under an assumed adequate recording, not measured on recordings: Full means that the information for a grade is in principle visible, not that a rater, still less an algorithm, can assign it reliably from routine clinical video. Studies of video-based administration offer a reference for clinician scoring: FMA and ARAT scores from standardised recordings agreed with direct observation (ICC 0.99) \cite{Amano2018RemoteFMAARAT}, BBS scores from patient-recorded videos correlated with face-to-face scores ($r$ = 0.95) \cite{Onal2025TeleBBS}, and telerehabilitation versions of the FMA-LE, BBS, POMA gait subscale and RMI showed moderate to good parallel-forms reliability (ICC 0.72--0.84) \cite{Zhang2024TeleMobilityScales}; these studies used trained raters and defined filming procedures and did not test automated extraction. Full codes assume that the task, support surfaces and any helper stay in frame, that distances have a visible reference, that hands are filmed close up for dexterity items and that the start signal is marked; a 30-m 6MWT walkway or centimetre-level reach distance may be hard to capture, and examiner perturbations in reactive balance items vary. For the 62 entries not fully available, Supplementary Data~1 specifies the additional source needed (counts overlap): patient interview or questionnaire (26 entries), audio for verbal cues or responses (22), the examiner's record of decisions such as the amount of assistance (22), manual examination of tone, reflexes, resistance or palpated events (16), records of usual performance (11) and force measurement (6). These judgements apply to RGB video only.

Fifth, the AI studies in Section~\ref{sec:framework} were selected during mapping to illustrate feature families, not identified systematically or appraised for quality, and cannot support conclusions about the prevalence, accuracy or validation status of approaches in the field.

\section{Conclusion}

Clinical rating scales remain the reference for post-stroke motor assessment, but their scores combine movement with clinical judgements that motion analysis does not capture by default. Across the scoring criteria of 57 tools, most score-defining information, including visible assistance and support, can in principle be obtained from an adequate recording, whereas the amount of help, verbal cues, supervision decisions, usual performance and content graded by force, tone, reflexes or patient report require other sources. This identifies which parts of clinical assessment motion analysis can target, not which scores can be automated. For AI-based assessment it implies four practices: quantify features anchored in scoring criteria and state whether each quantifies a criterion or only indicates a construct; annotate assistance events, the amount of help, verbal cues and safety decisions separately from movement; report feature definitions and recording protocols in enough detail for replication; and validate in stages, from analytical accuracy to construct validity and clinical utility, with established scales retained as clinical references.

\section*{Data Availability}
The tool library, all 196 mapping entries with coding rules, the second-reviewer check, the feature-family mapping and the guideline screening log are provided in Supplementary Table~\ref{tab:supplementary_tool_library_clinical}, Supplementary Table~\ref{tab:supplementary_tool_library_motion} and Supplementary Data~1.

\printbibliography

\input{Supplementary TableS1.tex}
\input{Supplementary Table S2.tex}

\noindent Supplementary Data~1 (Excel) contains the mapping entries underlying Supplementary Table~\ref{tab:supplementary_tool_library_motion} and Fig.~\ref{fig:video_availability}, the coding rules, the second-reviewer check, the mapping of computational feature families to clinical assessment content and the guideline full-text screening log.
\end{document}
